\documentclass[12pt]{article}
\usepackage[utf8]{inputenc}
\usepackage[T1]{fontenc}
\usepackage{hyperref}
\title{Pendau Language Resources for K-12 Curriculum Development}
\author{Compiled by Claude Sonnet 5 (claude-sonnet-5) for Andr\'as Kornai}
\date{}
\begin{document}
\maketitle
\section{Introduction}
Pendau (Ndau, Ndaoe, Umalasa; ISO 639-3 ums) is an Austronesian language of the Tomini-Tolitoli group in Central Sulawesi, Indonesia, with roughly 4,000 speakers. It is written in the Latin script. Children are schooled in Indonesian. The only substantial resource found is a scholarly grammar.
\section{Primary Educational Resources}
No primers, textbooks or school programmes for Pendau were found in search results.
\section{Linguistic and Reference Materials}
\begin{itemize}
\item Phil Quick, \emph{A grammar of the Pendau language of central Sulawesi, Indonesia} (Pacific Linguistics 590, 2007). \url{https://openresearch-repository.anu.edu.au/items/8aa0ec3c-bfaa-4990-ae51-006ba5dd30e6}
\item Quick, \emph{A grammar of the Pendau language} (ANU PhD thesis, 2003). \url{https://openresearch-repository.anu.edu.au/items/22f3f63e-f312-4db7-9cf1-d912589e26fe}
\item Glottolog reference record for Quick 2007. \url{https://glottolog.org/resource/reference/id/129448}
\item WorldCat record. \url{https://search.worldcat.org/oclc/314374225}
\end{itemize}
No standalone Pendau dictionary was found. Himmelmann's Tomini-Tolitoli word-list sourcebook (2001) is mentioned in the Endangered Languages entry.
\section{Religious and Community Resources}
None found.
\section{Digital and Audio Resources}
\begin{itemize}
\item Omniglot overview of Pendau. \url{https://omniglot.com/writing/pendau.htm}
\item Endangered Languages Project entry. \url{https://www.endangeredlanguages.com/lang/2769?hl=fr}
\end{itemize}
No audio or software was found.
\section{Curriculum Development Context}
The Quick grammar is the main basis for any orthography or curriculum work. Native-language teaching materials appear not to exist, so development would need local community partners, who were not identified in the searches.
\end{document}